\labsection{4}{Fuzzy Logic and ANFIS with Air-Quality Data}{sec:lab04-fuzzy}

\begin{labproject}{Combine expert rules with data-driven prediction}
\textbf{Scenario.} An environmental monitor needs an interpretable ventilation policy and a nonlinear predictor for carbon-monoxide concentration.\\
\textbf{Goal.} Implement a Mamdani-style fuzzy controller and a compact Sugeno ANFIS model in Python, then compare ANFIS with linear regression.

\begin{labmode}
Use \texttt{all\_datasets/lab04-AirQuality.csv} and run \texttt{all\_experiments/lab04\_fuzzy\_anfis\_worked.ipynb}. The notebook implements the fuzzy inference and ANFIS learning directly with NumPy/SciPy; no proprietary toolbox is required.
\end{labmode}

\begin{mlsteps}
  \item \textbf{Clean the data.} Parse the semicolon-separated CSV, replace \texttt{-200} with missing values, remove the empty tail, and report usable rows.
  \item \textbf{Build the fuzzy controller.} Use temperature and relative humidity as inputs and ventilation demand (0--100\%) as output. Define three membership functions per variable and clear expert rules.
  \item \textbf{Validate the rule base.} Evaluate five fixed operating points and inspect the control surface for smooth, sensible behaviour.
  \item \textbf{Prepare the ANFIS task.} Predict \texttt{CO(GT)} from \texttt{PT08.S1(CO)} and \texttt{PT08.S2(NMHC)} using a fixed 70/15/15 split with seed 42. Fit scaling on training data only.
  \item \textbf{Fit two model classes.} Train linear regression and ANFIS candidates with two and three membership functions per input. Select ANFIS size using validation RMSE only.
  \item \textbf{Evaluate once on test data.} Report RMSE and MAE, save predicted-versus-actual plots, and state whether ANFIS complexity is justified.
\end{mlsteps}

\begin{mltable}
\begin{tabularx}{\linewidth}{@{}>{\bfseries\leavevmode\color{MLNavy}}p{0.28\linewidth}XXX@{}}
\toprule
\rowcolor{MLPaleBlue!92}
Model & Validation RMSE & Test RMSE & Test MAE \\
\midrule
Linear baseline & & & \\
ANFIS (selected) & & & \\
\bottomrule
\end{tabularx}
\end{mltable}

\requiredoutput{Executed notebook, fuzzy membership/control-surface figures, split CSV, model metrics, and predicted-versus-actual figure.}
\end{labproject}

\begin{verification}
The test partition must remain untouched during ANFIS structure selection. Re-running the notebook with the saved seed must reproduce the split and reported metrics.
\end{verification}

\begin{labdeliverable}
Push the Python notebook, source, and results; create tag \texttt{lab04-final}. The README must explain the fuzzy rules, ANFIS inputs/target, split policy, baseline comparison, and one data limitation.
\end{labdeliverable}
